\section{Burst size, and why it is the quasi-stationary law}
\label{dist:sec:burst-size}

A cell presents two quantities to the outside world: when it ruptures, and how
much it releases. The second is settled first, because it is the law the
previous section has just computed.

\subsection{The burst-size law}
\label{dist:sec:burst-law}

The burst size $\Kb$ is the load immediately before rupture. For a burst in
state $k$ the process must be at $k$ and then rupture from there, which
happens at rate $\delta k$, so
\begin{equation}
\pr{\Kb = k} = \delta k \int_0^\infty p_k(t)\,\dd t.
\label{dist:eq:burstintegral}
\end{equation}
A discretisation makes this plain: over an interval of width $\Delta t$ in
which the process is at $k$, the probability that the burst occurs there is
$\pr{\Xt=k}\,\delta k\,\Delta t$, and summing over intervals as
$\Delta t\to0$ gives \cref{dist:eq:burstintegral}.

The integral is elementary once the sliding ratio is in view. With $P$ as in
\cref{dist:eq:pdef},
\begin{equation}
\dda{}{t}\Bigl[\frac{P(t)^k}{k\lambda}\Bigr]
=\frac{1}{\lambda}P(t)^{k-1}P'(t)
=\lrb{\frac{a-b}{b-aw}}^2 w\,P(t)^{k-1}
=p_k(t),
\label{dist:eq:antideriv}
\end{equation}
as one verifies from $P'(t)=\lambda(a-b)^2w/(b-aw)^2$. Since $P(0)=0$ and
$P(\infty)=1/a$,
\begin{equation}
\int_0^\infty p_k(t)\,\dd t=\frac{1}{k\lambda}\Bigl(\frac{1}{a}\Bigr)^{k},
\label{dist:eq:pkintegral}
\end{equation}
and \cref{dist:eq:burstintegral} gives the law.

\begin{theorem}[Burst-size distribution]
\label{dist:thm:burst}
Unconditionally, including the mass $b$ of realisations that never burst,
\begin{equation}
\pr{\Kb=k}=\frac{\delta}{\lambda}\,a^{-k},\qquad k\ge1;
\label{dist:eq:burstlaw}
\end{equation}
conditioned on the burst occurring,
\begin{equation}
\pr{\Kb=k\mid\text{burst}}=(a-1)\,a^{-k},\qquad k\ge1.
\label{dist:eq:burstlawcond}
\end{equation}
\end{theorem}

\begin{proof}
The unconditional form is the computation above. For the conditional form,
divide \cref{dist:eq:burstlaw} by the eventual-burst probability $1-b$:
\[
\pr{\Kb=k\mid\text{burst}}
=\frac{(\delta/\lambda)\,a^{-k}}{1-b}
=\frac{AB}{B}\,a^{-k}
=(a-1)\,a^{-k},
\]
using $\delta/\lambda=AB$ from \cref{dist:eq:ab}.
\end{proof}

Two consistency checks close the loop. Normalisation:
\[
\sum_{k\ge1}\frac{\delta}{\lambda}a^{-k}
=\frac{\delta}{\lambda}\,\frac{1}{a-1}
=\frac{AB}{A}=B=1-b,
\]
matching the eventual-burst probability. And the first moment:
\[
\sum_{k\ge1}k\,\frac{\delta}{\lambda}a^{-k}
=\frac{\delta}{\lambda}\,\frac{a}{(a-1)^2}
=\frac{aB}{A}=V_\infty,
\]
recovering \cref{dist:eq:vinf} --- which validates the burst-size law, the
formula for $V_\infty$ and the identity \cref{dist:eq:ab} at once.

\subsection{Why the two laws agree}
\label{dist:sec:identity}

Set \cref{dist:eq:burstlawcond} beside \cref{dist:eq:qsd}. They are the same
law.

\begin{theorem}[The burst size is the quasi-stationary law]
\label{dist:thm:identity}
The burst-size distribution conditioned on bursting is exactly the
quasi-stationary distribution of the load: both are geometric on
$\{1,2,\ldots\}$ with ratio $1/a$.
\end{theorem}

That the algebra agrees is not by itself an explanation, because the two
constructions share no step. On one side, condition the load on having avoided
both absorbing states and let $t\to\infty$: a limit along a single trajectory,
in which nothing is averaged and no rupture occurs. On the other, weight each
state by its rupture hazard $\delta k$ and integrate over all rupture times:
an average across trajectories, in which every realisation contributes at the
instant it dies and never afterwards. Nothing forces a conditional limit and
an occupation integral to agree, and for a killed process whose hazard was not
linear in the load they would not.
\Cref{dist:fig:two-mechanisms} sets the two side by side.

\begin{proof}
Read \cref{dist:eq:antideriv} as a change of clock rather than as an
antiderivative. Its $k=1$ case is
\begin{equation}
p_1(t)=\frac{P'(t)}{\lambda},
\label{dist:eq:ponelaw}
\end{equation}
and with $\sigma=P(t)$, which runs monotonically from $0$ to $1/a$ by
\cref{dist:thm:geometric}, the occupation density of state $k$ becomes a
perfect differential:
\begin{equation}
p_k(t)\,\dd t=\frac{1}{\lambda}\,\sigma^{k-1}\,\dd\sigma .
\label{dist:eq:clockchange}
\end{equation}
Hence
\begin{equation}
\pr{\Kb=k}=\delta k\!\int_0^\infty\! p_k\,\dd t
=\frac{\delta}{\lambda}\int_0^{1/a}\! k\sigma^{k-1}\,\dd\sigma
=\frac{\delta}{\lambda}\Bigl[\sigma^k\Bigr]_0^{1/a}
=\frac{\delta}{\lambda}\,a^{-k},
\label{dist:eq:telescope}
\end{equation}
which normalised on the bursting event is $(a-1)a^{-k}$, the law of
\cref{dist:eq:qsd}.
\end{proof}

The proof says more than the statement. The size-bias weight $k$ is exactly
the Jacobian that turns the occupation density of state $k$ into a perfect
differential, and integrating a perfect differential returns only its
endpoints. The burst-size law therefore depends on the whole trajectory
\emph{only} through the endpoint $\sigma_\infty=1/a$ that $P$ reaches, which
is by definition the quasi-stationary ratio. That is why the two laws
coincide: not a coincidence, and not a Yaglom limit, but size-biasing
telescoping the time integral onto the terminal value of $P$, the process's
own geometric ratio.

\begin{remark}[What the argument needs]
\label{dist:rem:criterion}
Nothing in the proof used the birth--death structure. The identity holds for
any killed process whose transient conditional law is geometric in a single
sliding ratio $P(t)$ satisfying $p_1=P'/\lambda$, with a killing hazard
linear in the load. Whether that criterion \emph{characterises} the class of
processes for which burst size equals the quasi-stationary law is not settled
here; see \cref{dist:sec:open-problems}.
\end{remark}

\begin{figure}[tbp]
\centering
\includegraphics[width=0.87\textwidth]{fig06_two_mechanisms/fig06}
\caption{Two mechanisms, one law. (a) The load conditioned on being neither
ruptured nor extinct, $p_n(t)/\Ifix(t)$, at $t=1,5,20$ from
\cref{dist:eq:stateprob,dist:eq:idihat}, with the quasi-stationary law
$(a-1)a^{-n}$ dashed; on the logarithmic scale each transient law is a
straight line whose slope rotates onto the limiting slope as $P(t)$ slides up
to $1/a$. (b) The burst-size law $(\delta/\lambda)a^{-k}$ of
\cref{dist:eq:burstlaw} (markers) and the same law conditioned on the burst
occurring (dashed) --- the identical line, reached without taking a limit at
all; the two series differ by the constant factor $1-b=0.812$. The first
panel arrives by conditioning on survival along a single trajectory, the
second by weighting each state by its rupture hazard and integrating across
trajectories. (c) The two routes on one axis, each computed on its own terms:
the conditional law of (a) evaluated at $t=60$, against the occupation
integral of (b) evaluated by numerical quadrature rather than in closed form.
They agree to $10^{-8}$, which is the accuracy of the quadrature grid and not
of the identity --- \cref{dist:thm:identity} makes the two exactly equal.
Parameters $(\lambda,\mu,\delta)=(1,0.2,0.05)$, so $a=1.0616$ and
$b=0.1884$; closed forms throughout except where panel (c) integrates
numerically, and no simulation anywhere.}
\label{dist:fig:two-mechanisms}
\end{figure}

\subsection{Size-biasing and the size of a late burst}
\label{dist:sec:size-bias}

The size-biasing that drives the telescoping has a reading of its own: it is
what an experiment sees when it times its observations.
A cell bursts at rate $\delta\Xt$, so given a burst at time $t$ its size is a
size-biased draw from $p_\cdot(t)$:
\begin{equation}
\pr{\Kb=k\mid\tau=t}=\frac{k\,p_k(t)}{J(t)},
\qquad
\E{\Kb\mid\tau=t}=\frac{K(t)}{J(t)}
=1+\frac{2\lambda}{\delta}\lrb{1-I(t)}.
\label{dist:eq:sizebias}
\end{equation}
Letting $t\to\infty$,
\begin{equation}
\E{\Kb\mid\tau=t}\xrightarrow[t\to\infty]{}
1+\frac{2\lambda}{\delta}(1-b)
=\frac{\langle X^2\rangle_{\rm QS}}{\langle X\rangle_{\rm QS}}
=\frac{a+1}{a-1},
\label{dist:eq:latelimit}
\end{equation}
the second equality from \cref{dist:eq:qsmoments}. A late burst is therefore
larger than an average one by the factor
\[
\frac{(a+1)/(a-1)}{a/(a-1)}=\frac{a+1}{a},
\]
which is $1.94$ at the working parameters but only $1.38$ at the
\textit{B.~anthracis} parameters of \cref{dist:sec:two-extremes}: the
often-quoted factor of two is the $a\to1$ limit, not a general fact.
\Cref{dist:fig:latemean} shows the rise.

The mechanism is a selection effect on the geometric law, not on individual
cells. Conditioning on a late rupture makes no cell grow; it discards the
cells that ruptured early, and early ruptures come disproportionately from
small loads, the hazard $\delta k$ being proportional to the load. What
survives is a sample already weighted towards large loads, which
\cref{dist:eq:sizebias} then weights once more at the moment of rupture.
\Cref{dist:fig:jointtauk} makes the tilt visible.

The size of a burst is now fully described. Its timing is not, and the two are
not independent --- which is what makes the remaining conditioning worth
doing.

\begin{figure}[tbp]
\centering
\includegraphics[width=0.62\textwidth]{fig07_late_burst_mean/fig07}
\caption{The size of a late burst at $(\lambda,\mu,\delta)=(1,0.2,0.05)$. The
conditional mean $\E{\Kb\mid\tau=t}=K(t)/J(t)$ of \cref{dist:eq:sizebias}
rises from $1$ at founding to $(a+1)/(a-1)=33.46$; the quasi-stationary mean
$a/(a-1)=17.23$ is marked for contrast, their ratio being $(a+1)/a=1.94$.}
\label{dist:fig:latemean}
\end{figure}

\begin{figure}[tbp]
\centering
\includegraphics[width=0.72\textwidth]{fig08_joint_tau_K/fig08}
\caption{Joint density, with respect to time and counting measure, of burst
time $\tau$ and burst size $\Kb$ for $(\lambda,\mu,\delta)=(1,0,0.1)$. The
overlay is the conditional mean $\E{\Kb\mid\tau=t}=K(t)/J(t)$, rising from
$1$ towards $1+2\lambda/\delta=21$; the upward tilt of the mass is
size-biasing made visible.}
\label{dist:fig:jointtauk}
\end{figure}
